% F22：本书原创矢量示意；依据所在节的定义、证明或明确模型。
\begin{figure}[htbp]
\centering
\begin{tikzpicture}[x=1cm,y=1cm,>=stealth,every node/.style={font=\small},line width=.65pt]
\node[] at (1.5,3.2) {\(j=2\)};
\fill[AnalysisBlue!10](0.0,0)rectangle(3.0,1.5);\draw[->,AnalysisBlue,thick](0.0,0)--(3.0,0);\draw[->,AnalysisBlue,thick](3.0,0)--(3.0,1.5);\draw[->,AnalysisBlue,thick](3.0,1.5)--(0.0,1.5);\draw[->,AnalysisBlue,thick](0.0,1.5)--(0.0,0);\node[] at (1.5,-0.5) {面积 \(1/2\)};
\node[] at (5.9,3.2) {\(j=4\)};
\fill[AnalysisBlue!10](4.4,0)rectangle(7.4,0.75);\draw[->,AnalysisBlue,thick](4.4,0)--(7.4,0);\draw[->,AnalysisBlue,thick](7.4,0)--(7.4,0.75);\draw[->,AnalysisBlue,thick](7.4,0.75)--(4.4,0.75);\draw[->,AnalysisBlue,thick](4.4,0.75)--(4.4,0);\node[] at (5.9,-0.5) {面积 \(1/4\)};
\node[] at (10.3,3.2) {\(j=8\)};
\fill[AnalysisBlue!10](8.8,0)rectangle(11.8,0.375);\draw[->,AnalysisBlue,thick](8.8,0)--(11.8,0);\draw[->,AnalysisBlue,thick](11.8,0)--(11.8,0.375);\draw[->,AnalysisBlue,thick](11.8,0.375)--(8.8,0.375);\draw[->,AnalysisBlue,thick](8.8,0.375)--(8.8,0);\node[] at (10.3,-0.5) {面积 \(1/8\)};
\node[] at (6,2.35) {\(Q_j=(0,1)\times(0,1/j),\qquad T_j=\partial[Q_j]\)};

\end{tikzpicture}
\caption[薄矩形边界的平坦收敛]{薄矩形边界的平坦收敛。长边反向靠拢。质量 \(\mathbf M(T_j)=2+2/j\) 趋于二，而填充面积给出 \(\mathbb F(T_j)\le1/j\)，故平坦极限为零；有向抵消不保存无向周长。}
\label{b1:fig:visual-f22}
\end{figure}
